\documentclass[12pt,a4paper]{article}
\usepackage[utf8]{inputenc}
\usepackage[T1]{fontenc}
\usepackage[spanish]{babel}
\usepackage{geometry}
\geometry{margin=2.5cm}
\usepackage{longtable}
\usepackage{booktabs}
\usepackage{graphicx}
\usepackage{hyperref}
\hypersetup{hidelinks}
\setlength{\parindent}{1.2em}
\setlength{\parskip}{0.35em}

% Completar estos campos con datos institucionales reales antes de entregar.
\newcommand{\Institucion}{[Nombre de la institución]}
\newcommand{\Programa}{[Programa o carrera]}
\newcommand{\Asignatura}{[Nombre de la asignatura]}
\newcommand{\Docente}{[Nombre del docente]}
\newcommand{\Integrantes}{[Nombres y matrículas de integrantes]}
\newcommand{\Grupo}{[Grupo]}
\newcommand{\LugarFecha}{[Ciudad y fecha de entrega]}

% Sustituir el recuadro por \includegraphics[width=.9\linewidth]{ruta/real.png}
% únicamente cuando exista una captura verificable y se haya anonimizado.
\newcommand{\capturaPendiente}[1]{%
  \fbox{\begin{minipage}[c][3.4cm][c]{0.86\linewidth}
    \centering\textbf{CAPTURA PENDIENTE}\\[0.5em]
    \small #1\\[0.4em]
    No representa una ejecución ni un resultado obtenido.
  \end{minipage}}%
}

\begin{document}
\begin{titlepage}
\centering
\IfFileExists{figuras/logo-institucional.png}{\includegraphics[height=2.5cm]{figuras/logo-institucional.png}\par\vspace{0.4cm}}{\fbox{\parbox[c][2cm][c]{5cm}{\centering Logotipo institucional pendiente}}\par\vspace{0.4cm}}
{\Large \Institucion\par}
\vspace{0.7cm}
{\large \Programa\par}
\vspace{1.3cm}
{\Large \Asignatura\par}
\vspace{1.8cm}
{\Huge\bfseries Práctica DevOps\par}
\vspace{0.6cm}
{\LARGE API REST, pruebas automatizadas y despliegue blue/green\par}
\vfill
\begin{tabular}{rl}
Integrantes: & \Integrantes \\
Docente: & \Docente \\
Grupo: & \Grupo \\
\end{tabular}
\vfill
{\LugarFecha\par}
\end{titlepage}

\tableofcontents
\clearpage

\section{Introducción}
Una aplicación de inventario no se limita a almacenar una lista de artículos. Para que la información sea útil, debe conservar relaciones entre usuarios, categorías, ubicaciones, departamentos, proveedores, marcas, unidades y proyectos, además de registrar préstamos, reservaciones y actividades de mantenimiento. La consulta de un recurso puede parecer una operación sencilla, pero el significado de una respuesta cambia cuando el identificador no existe, la solicitud carece de datos válidos o el almacenamiento no está disponible. Esta práctica propone una API REST desarrollada con Node.js 22 y Express para abordar esos casos de forma uniforme; SQLite actúa como almacenamiento embebido. La selección de esta base simplifica el montaje de una instancia académica, aunque obliga a prestar atención al archivo físico que contiene los datos y al acceso concurrente durante las actualizaciones de la aplicación.

El contrato solicitado comprende doce recursos con cinco operaciones por recurso: listar, consultar por identificador, crear, actualizar y eliminar. La combinación produce sesenta rutas CRUD bajo \texttt{/api}, a las que se suma un endpoint independiente \texttt{GET /api/health}. Esta distinción es relevante: comprobar que el proceso responde a \texttt{/api/health} sirve para decidir si una instancia puede recibir tráfico, pero no garantiza que todas las reglas de negocio o las sesenta rutas funcionen. El diseño de pruebas debe contrastar tanto las respuestas exitosas como las entradas inválidas, los identificadores ausentes y los errores de persistencia que se puedan reproducir sin comprometer datos de producción. Por ello se emplean Jest y Supertest sobre un entorno aislado; no debe confundirse la cantidad de pruebas con el porcentaje de código realmente ejercitado.

La calidad verificable también depende del proceso de integración. Se plantea que cada solicitud de cambio y cada envío a la rama principal active en GitHub Actions la instalación reproducible de dependencias y el comando de pruebas con cobertura. El criterio mínimo es setenta por ciento, medido con el reporte de Jest y exigido por la configuración del proyecto. Una ejecución que termine con error o que no calcule cobertura no satisface el criterio, aunque ciertas pruebas aisladas pasen. Solo después de una comprobación favorable en la rama principal corresponde publicar la imagen en Docker Hub: una etiqueta \texttt{latest} permite encontrar la versión más reciente y otra etiqueta correspondiente al SHA de GitHub permite identificar el artefacto que se intenta desplegar. Conservar esa identidad facilita investigar fallas y regresar a una versión anterior sin depender de una etiqueta mutable.

El empaquetado con Docker separa el código y sus dependencias del almacenamiento persistente. La aplicación debe escuchar en el puerto 3000 dentro del contenedor, mientras que el archivo de SQLite reside en \texttt{/app/data/database.sqlite} y se monta desde \texttt{/opt/devops-practica/data} en el host. Sin este montaje, cambiar o eliminar un contenedor podría borrar la información almacenada en su capa escribible. La ruta fija debe ir acompañada de permisos correctos, capacidad disponible y respaldos consistentes. Si SQLite opera en modo WAL, copiar solamente el archivo principal mientras hay escrituras puede omitir transacciones; además, dos versiones que compartan el archivo no deben iniciar migraciones incompatibles a la vez. Así, persistencia no equivale por sí sola a recuperación comprobada.

Para entregar el servicio se propone una instancia EC2 con Ubuntu, Docker y Nginx. Nginx escucha en el puerto público 80 y redirige las solicitudes a uno de dos contenedores publicados únicamente en el bucle local del host, mediante los puertos 3001 y 3002. En el mecanismo blue/green, el candidato se inicia en el puerto inactivo y se consulta su endpoint de salud antes de modificar el proxy. Si el candidato falla, el anterior debe continuar atendiendo. Si la verificación pasa, se valida la configuración de Nginx, se recarga y se vuelve a comprobar el acceso a través del puerto público antes de retirar el contenedor previo. Un error posterior a la conmutación requiere restaurar la ruta anterior; en el primer despliegue no existe versión previa. Estas condiciones hacen explícita la diferencia entre un simple reinicio y una sustitución controlada.

Finalmente, automatizar no elimina las responsabilidades operativas. Las credenciales de Docker Hub y la clave privada de SSH deben permanecer en secretos de GitHub; el host remoto se autentica mediante una clave conocida cuya huella se haya verificado, no mediante la desactivación de comprobaciones SSH. El grupo de seguridad limita el tráfico público a los puertos necesarios y los puertos internos no se exponen a Internet. Blue/green sobre una única máquina tampoco elimina el punto único de fallo de la instancia ni protege contra una migración irreversible del esquema. Por ese motivo, este documento distingue entre el procedimiento diseñado y la evidencia pendiente: las métricas, capturas, dirección pública y resultados solo se incorporarán cuando existan ejecuciones observables. La finalidad del reporte es permitir que otra persona compruebe el contrato, el umbral de calidad, la publicación y la continuidad del servicio sin atribuir al proyecto logros que no hayan sido medidos.
\clearpage

\section{Objetivo, alcance y arquitectura}
El objetivo es implementar y verificar una API REST de 60 rutas CRUD, publicarla como imagen Node.js 22 y desplegarla tras pasar una compuerta de pruebas y cobertura. La arquitectura prevista es: cliente HTTP $\rightarrow$ Nginx (EC2:80) $\rightarrow$ contenedor activo (host \texttt{127.0.0.1:3001} o \texttt{127.0.0.1:3002}, puerto interno 3000) $\rightarrow$ SQLite persistente. El flujo de entrega previsto es \texttt{pull\_request}/\texttt{push main} $\rightarrow$ GitHub Actions $\rightarrow$ Docker Hub $\rightarrow$ despliegue mediante \texttt{scripts/deploy.sh} en EC2. La especificación operativa y las instrucciones de preparación están en \texttt{README.md}; este texto no certifica que se hayan ejecutado.

\subsection{Matriz de endpoints}
Para cada recurso $R$ de la tabla se exige el mismo conjunto de cinco operaciones. La matriz siguiente define las 60 combinaciones sin omitir la variante con identificador; \texttt{GET /api/health} se cuenta aparte.

\begin{longtable}{p{0.27\textwidth}p{0.65\textwidth}}
\toprule
\textbf{Recurso $R$} & \textbf{Rutas HTTP requeridas} \\
\midrule
\endfirsthead
\toprule
\textbf{Recurso $R$} & \textbf{Rutas HTTP requeridas} \\
\midrule
\endhead
\texttt{items} & \multicolumn{1}{p{0.65\textwidth}}{\texttt{GET /api/R}, \texttt{GET /api/R/:id}, \texttt{POST /api/R}, \texttt{PUT /api/R/:id}, \texttt{DELETE /api/R/:id}} \\
\texttt{users} & Mismas cinco rutas, sustituyendo $R$ por \texttt{users}. \\
\texttt{categories} & Mismas cinco rutas, sustituyendo $R$ por \texttt{categories}. \\
\texttt{locations} & Mismas cinco rutas, sustituyendo $R$ por \texttt{locations}. \\
\texttt{departments} & Mismas cinco rutas, sustituyendo $R$ por \texttt{departments}. \\
\texttt{suppliers} & Mismas cinco rutas, sustituyendo $R$ por \texttt{suppliers}. \\
\texttt{brands} & Mismas cinco rutas, sustituyendo $R$ por \texttt{brands}. \\
\texttt{units} & Mismas cinco rutas, sustituyendo $R$ por \texttt{units}. \\
\texttt{projects} & Mismas cinco rutas, sustituyendo $R$ por \texttt{projects}. \\
\texttt{loans} & Mismas cinco rutas, sustituyendo $R$ por \texttt{loans}. \\
\texttt{reservations} & Mismas cinco rutas, sustituyendo $R$ por \texttt{reservations}. \\
\texttt{maintenance-records} & Mismas cinco rutas, sustituyendo $R$ por \texttt{maintenance-records}. \\
\bottomrule
\end{longtable}

\section{Metodología de pruebas e integración}
Se requiere ejecutar \texttt{npm ci}, \texttt{npm test} y \texttt{npm run test:coverage} con Node.js 22. Jest evalúa los casos y Supertest envía solicitudes a Express sin depender de una URL pública. Las pruebas deben emplear una SQLite aislada para impedir que una suite destruya los datos del host. Para cada familia CRUD conviene contrastar creación, lectura, actualización y eliminación, además de datos faltantes o inválidos e identificadores inexistentes; el código esperado debe reflejar el contrato implementado. El reporte de cobertura de Jest ha de imponer un \textbf{mínimo de 70\%} con un umbral que haga fallar el proceso cuando no se alcance. Se debe conservar el reporte real con porcentajes de líneas, ramas, funciones y sentencias que corresponda medir; no se deduce cobertura a partir del total de casos.

Según el modelo de flujos y eventos de GitHub Actions \cite{github}, el workflow \texttt{.github/workflows/main.yml} debe ejecutar pruebas y cobertura en PR y \texttt{push} a \texttt{main}. La publicación en Docker Hub y el despliegue se reservan al \texttt{push} exitoso a \texttt{main}. Los secretos requeridos son \texttt{DOCKERHUB\_USERNAME}, \texttt{DOCKERHUB\_TOKEN}, \texttt{EC2\_HOST}, \texttt{EC2\_USER}, \texttt{EC2\_SSH\_KEY} y \texttt{EC2\_KNOWN\_HOSTS}. La imagen publicada recibe las etiquetas \texttt{latest} y el SHA exacto del commit mediante las acciones oficiales de construcción y publicación de Docker \cite{docker}. La versión inactiva se prueba mediante \texttt{/api/health} antes de cambiar Nginx; ante error de salud o recarga, se conserva o restaura el puerto anterior. Se deben comparar el SHA de Actions, la etiqueta del registro y la imagen efectivamente ejecutada; las reglas de acceso a la instancia deben verificarse con la guía de EC2 \cite{aws}.

\section{Resultados locales y evidencia externa por completar}
Se ejecutaron localmente \texttt{npm test} y \texttt{npm run test:coverage} sobre la versión actual: \textbf{2 suites y 25 pruebas aprobadas}. La cobertura calculada de \texttt{index.js} fue \textbf{85,71\% de líneas}, 83,61\% de sentencias, 82,35\% de ramas y 83,33\% de funciones; las cuatro métricas superan el umbral de 70\% configurado en Jest. La prueba automatizada enumera 60 combinaciones de método y ruta CRUD distintas y un endpoint adicional de salud; los casos parametrizados ejercitan las cinco operaciones de cada uno de los doce recursos, además de validaciones, integridad referencial, persistencia y el socket heredado. También se construyó localmente la imagen con el Dockerfile; al ejecutar un contenedor con un volumen de datos se obtuvo HTTP 200 en \texttt{/api/health}, se creó un recurso y se constató que continuaba disponible después de detener y recrear el contenedor con el mismo volumen. Se trató de una prueba local con datos temporales, no de un despliegue EC2. Estas mediciones no prueban que GitHub Actions o Docker Hub hayan funcionado. Faltan capturas y enlaces para cotejarlas antes de entregar el PDF.

\begin{center}
\begin{tabular}{p{0.47\textwidth}p{0.42\textwidth}}
\toprule
\textbf{Comprobación} & \textbf{Resultado verificable por incorporar} \\
\midrule
Commit y fecha de la ejecución remota & [SHA y fecha reales: pendiente] \\
\texttt{npm test}; número de suites y casos & 2 suites y 25 pruebas aprobadas localmente \\
\texttt{npm run test:coverage}; umbral $\geq 70\%$ & 85,71\% líneas; 83,61\% sentencias; 82,35\% ramas; 83,33\% funciones. Aprobado localmente \\
60 rutas + \texttt{GET /api/health} & 61 combinaciones distintas; 12 casos CRUD parametrizados aprobados localmente \\
Docker local y persistencia & Imagen construida; salud HTTP 200; registro preservado después de reiniciar contenedor, comprobado localmente \\
PR y push a \texttt{main} en Actions & [Enlaces y estado de jobs: pendiente] \\
Etiquetas de Docker Hub & [Repositorio, \texttt{latest} y SHA: pendiente] \\
EC2, salud y rollback blue/green & [Comandos/salidas anonimizados: pendiente] \\
\bottomrule
\end{tabular}
\end{center}

Para sustituir cada recuadro numerado se debe tomar una captura auténtica en el entorno correspondiente, anonimizar direcciones o datos privados y cambiar \verb|\capturaPendiente{...}| por \verb|\includegraphics[width=.9\linewidth]{ruta/real.png}|. Conservar la leyenda y registrar fecha, commit y contexto en el texto. No usar imágenes simuladas como evidencia de ejecución.

La figura~\ref{fig:pruebas} deberá mostrar los comandos locales y el resumen legible de Jest. El procedimiento observado fue ejecutar primero la suite y luego repetirla con instrumentación de cobertura; en ambos casos pasaron las 25 pruebas. En la segunda ejecución Jest informó 85,71\% de líneas frente al mínimo de 70\%. La captura aún debe tomarse del terminal para que el lector pueda cotejar los números.

\begin{figure}[htbp]
\centering
\capturaPendiente{Insertar salida real de \texttt{npm test} y \texttt{npm run test:coverage} con umbral y porcentajes legibles.}
\caption{Evidencia de pruebas locales y cobertura (pendiente de captura real).}
\label{fig:pruebas}
\end{figure}

La figura~\ref{fig:ci} se obtendrá en la pestaña Actions tras abrir un PR y tras enviar un commit a \texttt{main}. Se comprobará que el PR ejecuta únicamente las pruebas y que el \texttt{push} no construye ni despliega si CI falla; su resultado permanece pendiente porque aún no existe una ejecución remota registrada en este reporte.

\begin{figure}[htbp]
\centering
\capturaPendiente{Insertar ejecución de GitHub Actions para PR y \texttt{push main}, con SHA y estados de CI y publicación visibles.}
\caption{Evidencia de CI/CD en GitHub Actions (pendiente de captura real).}
\label{fig:ci}
\end{figure}

Para la figura~\ref{fig:docker} se abrirá el repositorio de Docker Hub y se compararán sus etiquetas \texttt{latest} y SHA con el commit aprobado por Actions. Que dos nombres aparezcan en el registro no basta: se debe cotejar la identidad del artefacto desplegado. La publicación y esa comprobación externa todavía no se han realizado.

\begin{figure}[htbp]
\centering
\capturaPendiente{Insertar vista real del repositorio Docker Hub mostrando etiquetas \texttt{latest} y \texttt{<github.sha>}.}
\caption{Imagen publicada y etiquetas de Docker Hub (pendiente de captura real).}
\label{fig:docker}
\end{figure}

La figura~\ref{fig:despliegue} corresponderá a una consulta de salud desde un cliente externo y a la vista de los contenedores en EC2 antes y después de la conmutación. Se deberá registrar el SHA esperado en la respuesta de salud y repetir la consulta de un registro creado previamente para confirmar persistencia. Sin una EC2 configurada no es posible afirmar que se haya completado esa prueba.

\begin{figure}[htbp]
\centering
\capturaPendiente{Insertar comprobación real de \texttt{/api/health} mediante Nginx en puerto 80 y contenedor activo local; anonimizar host.}
\caption{Salud y conmutación blue/green en EC2 (pendiente de captura real).}
\label{fig:despliegue}
\end{figure}
\clearpage

\section{Limitaciones, riesgos y trabajo pendiente}
Dos contenedores pueden compartir el archivo SQLite durante el relevo, pero los escritores compiten y las migraciones de esquema incompatibles hacen inseguro el rollback de la aplicación. Es necesario preparar respaldos consistentes, ensayar su restauración y aplicar cambios de esquema compatibles con ambas versiones; \texttt{SQLITE\_BUSY} y WAL deben evaluarse según la carga y la versión de la biblioteca SQLite \cite{sqlite}. La disponibilidad depende de un solo host, de Nginx y de una recarga correcta: blue/green aquí reduce el riesgo de una actualización fallida, \emph{no} ofrece alta disponibilidad frente a la caída de EC2. El primer despliegue no tiene contenedor previo al que revertir.

Quedan pendientes las capturas de las pruebas locales, los enlaces de los jobs de Actions, las etiquetas publicadas y una demostración de salud y persistencia sobre EC2. La preparación externa y las reglas de red se describen en \texttt{README.md}; no se declara realizada ninguna acción en cuentas ajenas. La publicación de métricas sin salida verificable o de capturas con secretos incumpliría el criterio de evidencia responsable.

\section{Conclusión}
Este trabajo define un contrato medible para una API de doce recursos: cinco operaciones por cada uno de ellos y una comprobación de salud adicional. La combinación de pruebas HTTP, aislamiento de la base de datos y umbral mínimo de cobertura del 70\% proporciona una manera concreta de detectar regresiones antes de construir y publicar una imagen. La metodología evita interpretar el número de pruebas como prueba suficiente de calidad; la aceptación debe apoyarse en salidas reproducibles y en la revisión de los casos de error. La ejecución local de 25 pruebas y la medición de 85,71\% de líneas superaron el umbral, aunque todavía deben incorporarse capturas auténticas y verificarse de nuevo en el runner de GitHub.

La entrega blue/green propuesta separa el puerto público del proceso candidato: Nginx dirige tráfico a un contenedor sano mientras el otro se prepara en un puerto local, con una vía explícita de regreso si la conmutación falla. Esta medida disminuye el impacto potencial de una nueva versión, pero no resuelve la disponibilidad de un único servidor ni el posible cambio irreversible de datos compartidos. La siguiente etapa es completar las evidencias, verificar la identidad de la imagen desplegada, probar el respaldo/restauración de SQLite y ensayar un fallo controlado del candidato. Únicamente después de estas observaciones se podrá concluir si el sistema satisface los requisitos operativos y de cobertura en el ambiente de entrega.

\begin{thebibliography}{9}
\bibitem{github} GitHub Docs. \emph{GitHub Actions documentation}. \url{https://docs.github.com/en/actions}. Fecha de acceso: 6 de octubre de 2026.
\bibitem{docker} Docker Docs. \emph{Docker Build GitHub Actions}. \url{https://docs.docker.com/build/ci/github-actions/}. Fecha de acceso: 6 de octubre de 2026.
\bibitem{sqlite} SQLite. \emph{Write-Ahead Logging}. \url{https://www.sqlite.org/wal.html}. Fecha de acceso: 6 de octubre de 2026.
\bibitem{nginx} NGINX. \emph{Module ngx\_http\_proxy\_module}. \url{https://nginx.org/en/docs/http/ngx_http_proxy_module.html}. Fecha de acceso: 6 de octubre de 2026.
\bibitem{aws} Amazon Web Services. \emph{Authorizing access to an Amazon EC2 instance}. \url{https://docs.aws.amazon.com/AWSEC2/latest/UserGuide/authorizing-access-to-an-instance.html}. Fecha de acceso: 6 de octubre de 2026.
\end{thebibliography}
\end{document}
